\section*{Problems}

\subsection*{Problem statements}

% Remember
\begin{problema}
	\label{prob:ba2e465}
	Find the range, mean deviation, and standard deviation of the following ungrouped data arrays:
	\begin{enumerate}
		\item $12, 6, 7, 3, 15, 10, 18, 5$
		\item $9, 3, 8, 8, 9, 8, 9, 18.$
	\end{enumerate}
	Check your results with \texttt{Python}.
\end{problema}

% Understand
\begin{problema}
	\label{prob:c5833ad}
	If $d=X-P$ represents the deviations of a variable $X$ with respect to a constant pivot $P$, prove that the standard deviation is invariant under translations---that is, that its value does not depend on which pivot $P$ is chosen---and explain why this property is desirable for a measure of dispersion.
\end{problema}

% Apply
\begin{problema}
	\label{prob:d87f55c}
	Two groups of students have the following sample characteristics:
	\begin{itemize}
		\item Group A: $n_1 = 30$ students, $\bar{X}_1 = 75$ points, $s_1 = 8$ points
		\item Group B: $n_2 = 20$ students, $\bar{X}_2 = 75$ points, $s_2 = 12$ points
	\end{itemize}
	Find the standard deviation of the combined group of $n_1 + n_2 = 50$ students.
\end{problema}

% Analyze
\begin{problema}
	\label{prob:6112df0}
	Prove rigorously that for any data set (ungrouped or grouped):
	\begin{align}
		s &= \sqrt{\dfrac{\sum X^{2}}{N}-\left( \dfrac{\sum X}{N} \right)^{2}} = \sqrt{\overline{X^{2}}-\overline{X}^{2}}\\
		s &= \sqrt{\dfrac{\sum fX^{2}}{N}-\left( \dfrac{\sum fX}{N} \right)^{2}} = \sqrt{\overline{X^{2}}-\overline{X}^{2}}
	\end{align}
\end{problema}

% Evaluate
\begin{problema}
	\label{prob:dd18a2e}
	The coefficient of variation is defined as $CV = (s/\bar{X}) \times 100\%$. Compare and interpret relative variability in the following two practical situations:
	\begin{itemize}
		\item Salaries at two companies: Company A has $\bar{X}_A = \$45,000$ and $s_A = \$8,000$; Company B has $\bar{X}_B = \$38,000$ and $s_B = \$6,500$.
		\item Heights and weights in an anthropometric group: mean height $170$ cm with $s = 8$ cm; mean weight $70$ kg with $s = 10$ kg.
	\end{itemize}
	In each case, evaluate which of the two quantities being compared exhibits greater relative dispersion, and explain why the standard deviation alone (without dividing by the mean) would be insufficient for this comparison.
\end{problema}

% Create
\begin{problema}
	\label{prob:78de91c}
	A quality analyst groups the weight (in grams) of 50 manufactured pieces in the following table, with constant class width $c=4$ and pivot at the class mark of the central class $P=52$:
	\begin{center}
		\begin{tabular}{ccc}
			\hline
			Class mark $X_i$ & Frequency $f_i$ & $u_i=(X_i-P)/c$ \\
			\hline
			44 & 4  & $-2$ \\
			48 & 10 & $-1$ \\
			52 & 20 & $0$ \\
			56 & 12 & $1$ \\
			60 & 4  & $2$ \\
			\hline
		\end{tabular}
	\end{center}
	Design the complete solution using the change-of-scale-and-origin method ($s = c \cdot s_u$): compute $s_u$ from the $u_i$ column and use it to obtain the standard deviation $s$ of the piece weights in grams.
\end{problema}

\subsection*{Detailed solutions}

% Remember
\begin{solproblema}[prob:ba2e465]
	\begin{enumerate}
		\item For the array $12,6,7,3,15,10,18,5$:
		\begin{itemize}
			\item \textbf{Range:} $\max - \min = 18 - 3 = 15$
			\item \textbf{Mean:} $\bar{X} = \frac{12+6+7+3+15+10+18+5}{8} = \frac{76}{8} = 9.5$
			\item \textbf{Mean deviation ($MD$):}
			\begin{align}
				MD &= \frac{|12-9.5| + |6-9.5| + |7-9.5| + |3-9.5| + |15-9.5| + |10-9.5| + |18-9.5| + |5-9.5|}{8} \\
				&= \frac{2.5 + 3.5 + 2.5 + 6.5 + 5.5 + 0.5 + 8.5 + 4.5}{8} = \frac{34}{8} = 4.25
			\end{align}
			\item \textbf{Standard deviation ($s$):}
			\begin{align}
				s &= \sqrt{\frac{(12-9.5)^2 + (6-9.5)^2 + (7-9.5)^2 + (3-9.5)^2 + (15-9.5)^2 + (10-9.5)^2 + (18-9.5)^2 + (5-9.5)^2}{8}} \\
				&= \sqrt{\frac{6.25 + 12.25 + 6.25 + 42.25 + 30.25 + 0.25 + 72.25 + 20.25}{8}} = \sqrt{\frac{190}{8}} = \sqrt{23.75} \approx 4.87
			\end{align}
		\end{itemize}

		\item For the array $9,3,8,8,9,8,9,18$:
		\begin{itemize}
			\item \textbf{Range:} $18 - 3 = 15$
			\item \textbf{Mean:} $\bar{X} = \frac{9+3+8+8+9+8+9+18}{8} = \frac{72}{8} = 9$
			\item \textbf{Mean deviation ($MD$):}
			\begin{align}
				MD &= \frac{|9-9| + |3-9| + |8-9| + |8-9| + |9-9| + |8-9| + |9-9| + |18-9|}{8} \\
				&= \frac{0 + 6 + 1 + 1 + 0 + 1 + 0 + 9}{8} = \frac{18}{8} = 2.25
			\end{align}
			\item \textbf{Standard deviation ($s$):}
			\begin{align}
				s &= \sqrt{\frac{0^2 + (-6)^2 + (-1)^2 + (-1)^2 + 0^2 + (-1)^2 + 0^2 + 9^2}{8}} = \sqrt{\frac{120}{8}} = \sqrt{15} \approx 3.87
			\end{align}
		\end{itemize}
	\end{enumerate}
\end{solproblema}

\begin{lstlisting}[language=Python]
import numpy as np
# Data set 1
data1 = np.array([12, 6, 7, 3, 15, 10, 18, 5])
print("Data set 1:")
print(f"Range: {np.ptp(data1)}")
print(f"Mean: {np.mean(data1)}")
print(f"Mean deviation: {np.mean(np.abs(data1 - np.mean(data1)))}")
print(f"Standard deviation: {np.std(data1)}")
# Data set 2
data2 = np.array([9, 3, 8, 8, 9, 8, 9, 18])
print("\nData set 2:")
print(f"Range: {np.ptp(data2)}")
print(f"Mean: {np.mean(data2)}")
print(f"Mean deviation: {np.mean(np.abs(data2 - np.mean(data2)))}")
print(f"Standard deviation: {np.std(data2)}")
\end{lstlisting}

% Understand
\begin{solproblema}[prob:c5833ad]
	If $d_j = X_j - P$, then $X_j = d_j + P$. The mean of $X$ in terms of the pivot is:
	\begin{align}
		\bar{X} = \frac{\sum f_j X_j}{N} = \frac{\sum f_j (d_j + P)}{N} = \frac{\sum f_j d_j}{N} + P = \bar{d} + P.
	\end{align}
	The deviation from the mean transforms as:
	\begin{align}
		X_j - \bar{X} = (d_j + P) - (\bar{d} + P) = d_j - \bar{d}.
	\end{align}
	The term $P$ cancels exactly, so the variance is identical no matter which pivot is chosen:
	\begin{align}
		s^2 = \frac{\sum f_j (X_j - \bar{X})^2}{N} = \frac{\sum f_j (d_j - \bar{d})^2}{N}.
	\end{align}
	This invariance is desirable because the dispersion of a data set is an intrinsic property of how far the data are from each other, not of where an arbitrary measurement origin is placed: translating all data (changing the reference units, subtracting a constant) should not alter how much they vary among themselves.
\end{solproblema}

% Apply
\begin{solproblema}[prob:d87f55c]
	Since both groups have the same mean ($\bar{X}_1 = \bar{X}_2 = 75$), the combined mean is also $\bar{X} = 75$. In this case, with no difference between group means, the combined variance is the weighted average of the individual variances:
	\begin{align}
		s^2 &= \frac{n_1 s_1^2 + n_2 s_2^2}{n_1 + n_2} = \frac{30 \cdot 8^2 + 20 \cdot 12^2}{30 + 20} \\
		&= \frac{30 \cdot 64 + 20 \cdot 144}{50} = \frac{1920 + 2880}{50} = \frac{4800}{50} = 96
	\end{align}
	Therefore, the standard deviation of the combined group is:
	\begin{align}
		s = \sqrt{96} \approx 9.80 \text{ points.}
	\end{align}
\end{solproblema}

% Analyze
\begin{solproblema}[prob:6112df0]
	Start from the formal definition of variance (or population variance divided by $N$):
	\begin{align}
		s^2 &= \frac{\sum (X_j - \bar{X})^2}{N} = \frac{\sum (X_j^2 - 2X_j\bar{X} + \bar{X}^2)}{N} \\
		&= \frac{\sum X_j^2 - 2\bar{X}\sum X_j + N\bar{X}^2}{N} = \frac{\sum X_j^2}{N} - 2\bar{X}\frac{\sum X_j}{N} + \bar{X}^2 \\
		&= \frac{\sum X_j^2}{N} - 2\bar{X}\cdot\bar{X} + \bar{X}^2 = \frac{\sum X_j^2}{N} - \bar{X}^2 = \overline{X^2} - \overline{X}^2.
	\end{align}
	Taking the square root on both sides proves the identity:
	\begin{align}
		s = \sqrt{\overline{X^2} - \overline{X}^2}.
	\end{align}
	For data with frequencies, the same algebraic reasoning applies by replacing $\sum X_j^2$ with $\sum f_j X_j^2$ and $\sum X_j$ with $\sum f_j X_j$, since $N=\sum f_j$ and $\sum f_j$ acts as the total number of weighted observations.
\end{solproblema}

% Evaluate
\begin{solproblema}[prob:dd18a2e]
	\textbf{Salaries:}
	\begin{align}
		CV_A &= \frac{8000}{45000} \times 100\% \approx 17.78\% \\
		CV_B &= \frac{6500}{38000} \times 100\% \approx 17.11\%
	\end{align}
	Both companies show almost similar relative salary dispersion (around 17\%), although Company A has slightly greater variability relative to its average payroll amount. Directly comparing $s_A=\$8{,}000$ against $s_B=\$6{,}500$ would be misleading, because it ignores that Company A also pays higher average salaries.

	\textbf{Heights vs. weights:}
	\begin{align}
		CV_{\text{height}} &= \frac{8}{170} \times 100\% \approx 4.71\% \\
		CV_{\text{weight}} &= \frac{10}{70} \times 100\% \approx 14.29\%
	\end{align}
	Although the physical units are different (cm vs.\ kg), the coefficient of variation enables a direct dimensionless comparison. The conclusion is that weight has roughly triple the relative variability of height in this population. In both cases, the standard deviation alone is insufficient for comparison because it mixes the absolute scale of the variable with its dispersion: two variables with the same $s$ but very different means (or variables measured in different units, such as cm and kg) are not comparable in terms of relative variability without normalizing by the mean.
\end{solproblema}

% Create
\begin{solproblema}[prob:78de91c]
	With $N=\sum f_i = 4+10+20+12+4=50$, compute $\sum f_iu_i$ and $\sum f_iu_i^2$:
	\begin{align}
		\sum f_iu_i &= 4(-2)+10(-1)+20(0)+12(1)+4(2) = -8-10+0+12+8 = 2 \\
		\sum f_iu_i^2 &= 4(4)+10(1)+20(0)+12(1)+4(4) = 16+10+0+12+16 = 54
	\end{align}
	The coded standard deviation is:
	\begin{align}
		s_u = \sqrt{\dfrac{\sum f_iu_i^2}{N} - \left(\dfrac{\sum f_iu_i}{N}\right)^2} = \sqrt{\dfrac{54}{50} - \left(\dfrac{2}{50}\right)^2} = \sqrt{1.08 - 0.0016} = \sqrt{1.0784} \approx 1.0385.
	\end{align}
	Applying the scale change $s = c \cdot s_u$ with $c=4$:
	\begin{align}
		s = 4 \times 1.0385 \approx 4.15 \text{ grams.}
	\end{align}
	The coded method avoids handling the larger class-mark values directly (44 to 60), operating instead with small integers ($-2$ to $2$).
\end{solproblema}
